
\documentclass[11pt]{article}
\usepackage{fullpage}
\usepackage[left=1in,top=1in,right=1in,bottom=1in,headheight=3ex,headsep=3ex]{geometry}
\usepackage{graphicx}
\usepackage{float}
\usepackage{enumerate}
\usepackage{sgamevar, tikz} % Game theory packages
\usepackage{pgfplots}

\usepackage{sgamevar, tikz} % Game theory packages
\usetikzlibrary{trees, calc} % For extensive form games

\newcommand{\blankline}{\quad\pagebreak[2]}
%%%%%%%%%%%%%%%%%%%%%%%%%%%%%%%%%%%%%%%%%%%%%%%%%%%%%%%%%%%%%%

% Modify Course title, instructor name, semester here %%%%%%%%

\title{Microeconom\'ia Aplicada (MAE)}
\author{Ayudantia 1: Juegos Est\'aticos con informaci\'on completa \\ Profesor: Jose Manuel Paz y Mi\~no \\ Ayudante: Vicente Merrill}
\date{}

% Don't touch this %%%%%%%%%%%%%%%%%%%%%%%%%%%%%%%%%%%%%%%%%%%
\usepackage[sc]{mathpazo}
\linespread{1.05} % Palatino needs more leading (space between lines)
\usepackage[T1]{fontenc}
\usepackage{setspace}
\usepackage{multicol}
\usepackage{fancyhdr,lastpage}
\usepackage{url}
\pagestyle{fancy}
\usepackage{hyperref}
\usepackage{lastpage}
\usepackage{amsmath}
\usepackage{layout}

\lhead{}
\chead{}
%%%%%%%%%%%%%%%%%%%%%%%%%%%%%%%%%%%%%%%%%%%%%%%%%%%%%%%%%%%%%%

% Modify header here %%%%%%%%%%%%%%%%%%%%%%%%%%%%%%%%%%%%%%%%%
\rhead{\footnotesize Microeconom\'ia Aplicada, Oto\~no 2025}

%%%%%%%%%%%%%%%%%%%%%%%%%%%%%%%%%%%%%%%%%%%%%%%%%%%%%%%%%%%%%%
% Don't touch this %%%%%%%%%%%%%%%%%%%%%%%%%%%%%%%%%%%%%%%%%%%
\lfoot{}
\cfoot{\small \thepage/\pageref*{LastPage}}
\rfoot{}

% \usepackage{array, xcolor}
% \usepackage{color,hyperref}
% \definecolor{clemsonorange}{HTML}{EA6A20}
% \hypersetup{colorlinks,breaklinks,linkcolor=clemsonorange,urlcolor=clemsonorange,anchorcolor=clemsonorange,citecolor=black}

%===================Startup========================
\begin{document}

\maketitle

\blankline

\setlength{\parindent}{0pt}.

1.- (Estrategias mixtas) Asuma tiene el siguiente juego. Se sabe que para el jugador 1, $p_{x}$ es la probabilidad de jugar X, $p_{y}$ es la probabilidad de jugar Y y $1-p_{x}-p_{y}$ es la probabilidad de jugar Z. Para el jugador 2, $q$ es la probabilidad de jugar X y $1-q$ es la probabilidad de jugar Y.

\begin{center}
	 \begin{game}{3}{2}[\textbf{1}][\textbf{2}]
                  \>$X$         \>$Y$\\
        $X$       \>$5,5$       \>$8,8$\\
        $Y$       \>$9,9$       \>$2,2$\\
        $Z$       \>$7,7$       \>$6,6$
 \end{game}
\end{center}

\begin{enumerate}
	\item Asuma $q=0.5$, ¿cu\'al es la mejor respuesta del jugador 1?
	\item Asuma $p_{x} = 0.2$ y $p_{y} = 0.2$, ¿cu\'al es la mejor respuesta del jugador 2?
	\item Encuentre todos los equilibrios de Nash (en estrategias puras y mixtas). 
	\item De acuerdo a su respuesta anterior, calcule la probabilidad de cada ocurrencia de cada resultado para el equilibrio de Nash con estrategias mixtas.
\end{enumerate}

2.- (Un modelo de lobbying) Se tienen 2 grupos de interes (A y B), que buscan influenciar a una pol\'itica de gobierno. Cada uno tiene una preferencia de pol\'itica distinta. A prefiere pol\'itica 0 y B prefiere pol\'itica 1. Cada grupo decide la contribuci\'on monetaria a oficiales $s_{i} \in [0,1]$ para $i=A,B$. La pol\'itica de gobierno esta influenciada de la siguiente manera
\begin{align*}
	x(s_{A},s_{B}) = 0.5 - s_{A} + s_{B}
\end{align*} por lo que si ning\'un grupo contribuye se ejecuta la politica $0.5$. Finalmente, los grupos de interes tienen la siguiente funci\'on de utilidad
\begin{align*}
	u_{A}(s_{A},s_{B}) &= -(x(s_{A},s_{B}))^{2} - s_{A} \\
	u_{B}(s_{A},s_{B}) &= -(1 - x(s_{A},s_{B}))^{2} - s_{B}
\end{align*}

\begin{enumerate}
	\item Encuentre el/los equilibrios de Nash.
	\item ¿Existe alg\'un equilibrio de Nash que Pareto domine al resto?
\end{enumerate}

3.- (Describiendo el espacio de utilidades posibles) Asuma el siguiente juego.
\begin{center}
	 \begin{game}{2}{2}[\textbf{1}][\textbf{2}]
                  \>$L$         \>$R$      \\
        $U$       \>$6,6$       \>$2,7$    \\
        $D$       \>$7,2$       \>$0,0$     
 \end{game}
 \end{center}

\begin{enumerate}
	\item Encuentre todos los equilibrios de Nash (en estrategias puras y mixtas) y grafique el convex hull de equilibrios de Nash.
	\item Caracterice las condiciones que debe cumplir un equilibrio correlacionado. ¿Es posible encontrar un equilibrio correlacionado fuera del convex hull de equilibrios de Nash?
\end{enumerate}

\end{document}  